\documentclass[11pt]{article}

\usepackage[T1]{fontenc}
\usepackage{lmodern}
\usepackage[margin=1in]{geometry}
\usepackage{amsmath,amssymb,amsthm,mathtools}
\usepackage{booktabs,longtable,array,tabularx}
\usepackage{microtype}
\usepackage[hidelinks]{hyperref}

\newtheorem{theorem}{Theorem}[section]
\newtheorem{lemma}[theorem]{Lemma}
\newtheorem{proposition}[theorem]{Proposition}
\newtheorem{corollary}[theorem]{Corollary}
\theoremstyle{definition}
\newtheorem{definition}[theorem]{Definition}
\theoremstyle{remark}
\newtheorem{remark}[theorem]{Remark}

\newcommand{\Q}{\mathbb Q}
\newcommand{\N}{\mathbb N}
\newcommand{\strictbr}[1]{\mathopen{[}#1\mathclose{]}_{<}}
\newcommand{\rhoc}{\rho_{c}}
\newcommand{\rhozero}{\rho_{0}}
\newcommand{\frB}{\mathfrak B}
\newcommand{\frC}{\mathfrak C}
\newcommand{\frZ}{\mathfrak Z}

\title{Exact-ledger decomputerization:\\Packets A and B}
\author{}
\date{16 July 2026}

\begin{document}
\maketitle

\begin{abstract}
This note replaces two groups of executable exact-arithmetic checks in the
three-dimensional spherical-shell P\'olya proof.  Packet~A proves once and
for all the strict-bracket convention, odd-multiplicity block sums, the
optical product-deficit identity, and monotonicity of propagated zero and
channel cuts.  Packet~B gives exact hand-checkable tables for the five
moving \(3z\)-crossings and for the seven first-excluded checkpoint
channels.  It then proves, without a finite search, that the five printed
checkpoint inventories are exhaustive.  Every numerical comparison below
is an equality of rational numbers, a cross multiplication, or a
positive-side squaring.  No program output, decimal approximation, or
interval calculation is a premise.
\end{abstract}

\tableofcontents

\section{Scope and notation}

Throughout \(0<\rho<1\).  Put
\[
 z=z_\rho:=\frac{\pi}{1-\rho},\qquad
 k_m(\rho)^2:=z_\rho^2+m(m+1),
\]
and define the two ratio endpoints
\[
 \rhozero:=\frac1{\sqrt{337}},
 \qquad
 \rhoc:=\frac1{1+2\pi}.
\]
The high-checkpoint argument uses the strict proxy
\begin{equation}
 b_{\ell,n}(z)^2
 :=\max\{n^2z^2+\ell(\ell+1),\,\beta_{\ell,n}^2\}.
 \label{eq:strict-proxy}
\end{equation}
A channel may be counted at \(K\) only when
\(b_{\ell,n}(z)<K\).  The constants \(\beta_{\ell,n}\) are strict lower
bounds for the corresponding ball zeros.  The zero bounds used in the
checkpoint table are displayed explicitly in Section~\ref{sec:checkpoint}.

The purpose of this note is deliberately limited.  It replaces the
general exact-arithmetic and checkpoint-ordering rows of the old ledger; it
does not replace the zero-sign proofs, high-event payment registry, seam
arithmetic, optical endpoint screens, or formal owner subtraction.

\section{A rational enclosure for \texorpdfstring{\(\pi\)}{pi}}
\label{sec:pi}

\begin{lemma}[Hand enclosure for \(\pi\)]
\label{lem:pi}
One has
\begin{equation}
 \frac{333}{106}<\pi<\frac{22}{7},
 \qquad
 \frac1\pi<\frac{106}{333}.
 \label{eq:pi-bounds}
\end{equation}
\end{lemma}

\begin{proof}
For \(0<x\leq1\), the alternating series for \(\arctan x\) gives a
strict lower bound after a negative term and a strict upper bound after a
positive term.  Hence
\[
 \arctan\frac15>
 \frac15-\frac1{3\cdot5^3}+\frac1{5\cdot5^5}
 -\frac1{7\cdot5^7}
 =\frac{323852}{1640625},
 \qquad
 \arctan\frac1{239}<\frac1{239}.
\]
Machin's identity therefore yields
\[
 \pi=16\arctan\frac15-4\arctan\frac1{239}
 >\frac{1231847548}{392109375},
\]
and direct subtraction gives the positive margin
\[
 \frac{1231847548}{392109375}-\frac{333}{106}
 =\frac{3418213}{41563593750}>0.
\]
For the upper bound, polynomial division and elementary integration give
\[
 0<\int_0^1\frac{x^4(1-x)^4}{1+x^2}\,dx
 =\frac{22}{7}-\pi.
\]
Finally, reciprocation of the positive lower bound for \(\pi\) gives the
last inequality.  Notice also that
\(22/7-333/106=1/742>0\), so the enclosure is nonempty.
\end{proof}

\section{Packet A: universal exact-arithmetic lemmas}

\subsection{Strict floors and equality owners}

\begin{definition}
For \(x\in\mathbb R\), let
\[
 \strictbr{x}:=\#\{q\in\mathbb Z_{\geq1}:q<x\}.
\]
Thus the bracket counts positive integers under a \emph{strict} cutoff.
\end{definition}

\begin{lemma}[Strict-bracket identity]
\label{lem:strict-bracket}
For every real \(x\),
\begin{equation}
 \strictbr{x}=\max\{0,\lceil x\rceil-1\}.
 \label{eq:strict-bracket}
\end{equation}
In particular, for every integer \(m\geq1\),
\[
 \strictbr{m}=m-1,
 \qquad
 \strictbr{m+\tfrac12}=m.
\]
Consequently an integer equality wall is owned by the excluded side.
\end{lemma}

\begin{proof}
If \(x\leq1\), both sides of \eqref{eq:strict-bracket} vanish.  If
\(x>1\), the positive integers below \(x\) are exactly
\(1,\ldots,\lceil x\rceil-1\).  This also proves the two specializations.
\end{proof}

The same lemma handles the two equality conventions used later.  If
\(A+1/4=n\in\mathbb Z_{\geq1}\), then
\(\strictbr{A+1/4}=n-1\).  If \(a=m\pi\) in the optical radial
decomposition, the strict radial layers are represented by
\((N,t)=(m-1,1)\), not by \((m,0)\).

\subsection{Odd multiplicities and block caps}

\begin{lemma}[Odd-sum and block-cap identity]
\label{lem:block-cap}
For every integer \(h\geq0\),
\begin{equation}
 \sum_{\ell=0}^{h}(2\ell+1)=(h+1)^2,
 \qquad
 \sum_{\ell=0}^{h-1}(2\ell+1)=h^2.
 \label{eq:odd-sum}
\end{equation}
More generally, if radial layer \(n\) contains a complete angular block
\(0\leq\ell\leq h_n\), with \(h_n=-1\) denoting an empty block, then its
total multiplicity is
\begin{equation}
 \sum_n (h_n+1)^2.
 \label{eq:block-cap}
\end{equation}
If \(x^2=h(h+1)\) is a strict product equality wall, the equality channel
\(\ell=h\) is absent and the contribution is \(h^2\).
\end{lemma}

\begin{proof}
The first identity follows either by induction or by
\((\ell+1)^2-\ell^2=2\ell+1\).  Summing over independent radial layers
gives \eqref{eq:block-cap}.  At \(x^2=h(h+1)\),
\(\sqrt{x^2+1/4}-1/2=h\); Lemma~\ref{lem:strict-bracket} excludes
\(\ell=h\), leaving \(0\leq\ell<h\) and hence \(h^2\).
\end{proof}

For completeness, the cap arithmetic collapsed by this lemma is printed
in compact form below.  A dash means that another analytic argument has
declared the cell impossible; the lemma asserts only the arithmetic of
the non-dash entries.
\[
\begin{array}{c|ccccc}
k_8:H\backslash h&\varnothing&0&1&2&3\\ \hline
7&64&65&68&73&80\\
6&49&50&53&-&-\\
5&36&37&40&45&52\\
4&25&26&29&34&41
\end{array}
\qquad
\begin{array}{c|cccccc}
k_9:H\backslash h&\varnothing&0&1&2&3&4\\ \hline
8&81&-&-&-&-&-\\
7&64&65&68&73&-&-\\
6&49&50&53&58&65&74\\
5&36&37&40&45&52&61\\
4&25&26&29&34&41&50
\end{array}
\]
Indeed every entry is \((H+1)^2\) plus \(0\) or \((h+1)^2\).
At \(k_7\), the first blocks \(H=4,5,6\) give
\(25,36,49\), and second channels \(0,1\) contribute \(1+3=4\).
The remaining printed caps decompose as
\begin{align*}
k_{10}:\quad&10^2;\quad9^2,\ 9^2+4^2;\\
&8^2,\ 8^2+4^2,\ 8^2+5^2,\ 8^2+6^2;\\
&7^2+5^2,\ 7^2+4^2+2^2,\ 7^2+5^2+2^2,
  7^2+6^2,\ 7^2+6^2+2^2\\
&=100;\quad81,97;\quad64,80,89,100;\quad74,69,78,85,89,\\[2mm]
k_{11}:\quad&11^2;\quad10^2,\ 10^2+5^2;\quad
9^2+5^2,\ 9^2+5^2+2^2;\quad8^2+5^2+2^2\\
&=121;\quad100,125;\quad106,110;\quad93.
\end{align*}
The conservative cap-74 row is the same identity \(7^2+5^2=74\).

The six lower inventories similarly give, in their printed order,
\[
 (6^2+3^2,\ 6^2+3^2+1,\ 7^2+3^2+1,
  7^2+4^2+1,\ 7^2+4^2+1+3,\ 7^2+5^2+1+3)
 =(45,46,59,66,69,78).
\]
Finally, the ten lower-\(d\) caps are the successive odd sums
\[
\begin{array}{c|c@{\qquad}c|c}
\text{channels}&\text{cap}&\text{channels}&\text{cap}\\ \hline
1&1&1+3&4\\
1+3+5&9&1+3+5+1&10\\
1+3+5+7&16&1+3+5+7+1&17\\
1+3+5+7+9+1&26&1+3+5+7+9+1+3&29\\
1+3+5+7+9+11+1+3&40&
1+3+5+7+9+11+1+3+5&45.
\end{array}
\]

\subsection{The symbolic optical product deficit}

\begin{proposition}[Product-deficit identity and uniform reserve]
\label{prop:product-deficit}
Let \(N\geq1\) be an integer, \(0<t\leq1\), and \(x=N+t\).  Then
\begin{align}
 \frac23x^3-Nx^2+\sum_{n=1}^N n^2
 &=\frac{N^2}{2}+N\left(t^2+\frac16\right)+\frac{2t^3}{3}.
 \label{eq:product-identity}
\end{align}
For
\(C_D=1382/3125\), the right-hand side is strictly larger than
\(C_D(N+t)^2\).  More precisely, its difference from
\(C_D(N+t)^2\) is at least
\begin{equation}
 \frac{1822532}{91552734375}>0.
 \label{eq:product-margin}
\end{equation}
\end{proposition}

\begin{proof}
Insert
\(\sum_{n=1}^Nn^2=N(N+1)(2N+1)/6\) and expand; this proves
\eqref{eq:product-identity} for all \(N,t\), rather than at sampled
values.  Put \(C=C_D\), \(A=1/2-C=361/6250\), and let \(G_N(t)\)
denote the difference under consideration.  Then
\[
 G_N(t)=AN^2+N\left(t^2-2Ct+\frac16\right)
       +\frac{2t^3}{3}-Ct^2.
\]
Consequently
\[
 G_{N+1}(t)-G_N(t)
 =A(2N+1)+t^2-2Ct+\frac16.
\]
Since \(0<C<1\), the last quadratic has minimum \(1/6-C^2\) on
\((0,1]\).  For \(N\geq1\),
\[
 G_{N+1}(t)-G_N(t)
 \geq3A+\frac16-C^2
 =\frac{4229603}{29296875}>0.
\]
It remains to minimize \(G_1\).  Direct differentiation gives
\[
 G_1'(t)=2(t-C)(t+1),
\]
so its unique minimum on \(0<t\leq1\) is at \(t=C\).  Exact
substitution gives
\[
 G_1(C)=\frac{1822532}{91552734375}>0,
\]
which proves both assertions.
\end{proof}

\subsection{Monotonic propagated cuts}

\begin{lemma}[Zero-cut/channel-cut monotonicity]
\label{lem:cut-monotonicity}
Fix \(n\geq2\), a checkpoint \(m\), and \(M=m(m+1)\).  Let
\(\mathcal A_n\) be a finite collection of anchor pairs
\((j,q_{j,n}^2)\), where
\(j_{j+1/2,n}^2>q_{j,n}^2\) at each anchor.  For every \(\ell\)
above at least one anchor, set
\begin{align}
 \frB_n(\ell)
 &:=\max_{(j,q_{j,n}^2)\in\mathcal A_n,\ j\leq\ell}
 \bigl(q_{j,n}^2+\ell(\ell+1)-j(j+1)\bigr),\label{eq:propagated-B}\\
 \frZ_{m,n}(\ell)&:=\frB_n(\ell)-M,\label{eq:zero-cut}\\
 \frC_{m,n}(\ell)&:=
 \frac{M-\ell(\ell+1)}{n^2-1}.\label{eq:channel-cut}
\end{align}
Then
\begin{align}
 \frZ_{m,n}(\ell+1)&\geq
 \frZ_{m,n}(\ell)+2(\ell+1)>\frZ_{m,n}(\ell),\label{eq:Z-monotone}\\
 \frC_{m,n}(\ell+1)&=
 \frC_{m,n}(\ell)-\frac{2(\ell+1)}{n^2-1}
 <\frC_{m,n}(\ell).\label{eq:C-monotone}
\end{align}
In particular, if
\(\frZ_{m,n}(\ell_0)\geq\frC_{m,n}(\ell_0)\), then no channel
\((\ell,n)\) with \(\ell\geq\ell_0\) can satisfy the strict proxy
\eqref{eq:strict-proxy} at \(K=k_m\).
\end{lemma}

\begin{proof}
Every old candidate in \eqref{eq:propagated-B} gains exactly
\((\ell+1)(\ell+2)-\ell(\ell+1)=2(\ell+1)\) when \(\ell\) is
replaced by \(\ell+1\); a new anchor can only increase the maximum.  This
proves \eqref{eq:Z-monotone}.  Formula \eqref{eq:C-monotone} is direct.

If a channel were counted at \(k_m\), its product part would give
\[
 (n^2-1)z^2<M-\ell(\ell+1),
 \quad\text{hence}\quad z^2<\frC_{m,n}(\ell),
\]
whereas its propagated ball-zero part would give
\[
 \frB_n(\ell)<z^2+M,
 \quad\text{hence}\quad z^2>\frZ_{m,n}(\ell).
\]
These strict inequalities are incompatible when
\(\frZ\geq\frC\).  Equations \eqref{eq:Z-monotone} and
\eqref{eq:C-monotone} propagate the incompatibility to every larger
angular index.
\end{proof}

\section{Packet B: moving-\texorpdfstring{\(3z\)}{3z} ordering}
\label{sec:moving}

We first record rational radical enclosures.  All six follow by positive
squaring; the exact square margins are
\[
\begin{array}{c|c@{\qquad}c|c}
\text{lower comparison}&\text{margin}&\text{upper comparison}&\text{margin}\\ \hline
(183/10)^2<337&211/100&337<(92/5)^2&39/25\\
(1051/100)^2<7073/64&2221/40000&7073/64<(1052/100)^2&6191/40000\\
(1067/100)^2<114&1511/10000&114<(107/10)^2&49/100.
\end{array}
\]
Thus
\begin{equation}
 \frac{183}{10}<\sqrt{337}<\frac{92}{5},\quad
 \frac{1051}{100}<\frac{\sqrt{7073}}8<\frac{1052}{100},\quad
 \frac{1067}{100}<\sqrt{114}<\frac{107}{10}.
 \label{eq:radical-bounds}
\end{equation}

\begin{lemma}[Endpoint and fixed-wall order]
\label{lem:fixed-order}
With \(d=\sqrt{337}/2\),
\begin{equation}
 d<3z_{\rhozero}<10<\frac{81}{8}<\frac{21}{2}
 <\frac{\sqrt{7073}}8<\sqrt{114}<3z_{\rhoc}.
 \label{eq:full-fixed-order}
\end{equation}
Moreover
\((81/8)^2+8=7073/64\).
\end{lemma}

\begin{proof}
The nontrivial endpoint estimates have the following rational margins.
From \eqref{eq:pi-bounds} and \eqref{eq:radical-bounds},
\[
 3z_{\rhozero}>3\pi>\frac{999}{106},
 \qquad d<\frac{46}{5},
 \qquad
 \frac{999}{106}-\frac{46}{5}=\frac{119}{530}>0.
\]
Also \(\rhozero<10/183\), whence
\[
 3z_{\rhozero}
 <\frac{66}{7}\frac{183}{173}
 =\frac{12078}{1211},
 \qquad
 10-\frac{12078}{1211}=\frac{32}{1211}>0.
\]
At the other endpoint, \(z_{\rhoc}=\pi+1/2\); therefore
\[
 3z_{\rhoc}>\frac{579}{53},
 \qquad
 \sqrt{114}<\frac{107}{10},
 \qquad
 \frac{579}{53}-\frac{107}{10}=\frac{119}{530}>0.
\]
The rational middle comparisons are immediate.  For the two radical
steps one may square positive quantities:
\[
 \frac{7073}{64}-\left(\frac{21}{2}\right)^2=\frac{17}{64}>0,
 \qquad
 114-\frac{7073}{64}=\frac{223}{64}>0.
\]
Finally,
\((81/8)^2+8=(6561+512)/64=7073/64\).
\end{proof}

For a positive wall \(w\), the unique solution of \(3z_\rho=w\) is
\[
 r_w:=1-\frac{3\pi}{w}.
\]
The following table encloses each crossing between explicit rational
numbers.  The bounds use only Lemma~\ref{lem:pi} and
\eqref{eq:radical-bounds}; for instance \(r_w\) increases with \(w\) and
decreases with \(\pi\).

\begin{center}
\renewcommand{\arraystretch}{1.25}
\begin{tabular}{c|c|c}
wall \(w\)&exact rational enclosure for \(w\)&exact enclosure for \(r_w\)\\ \hline
\(10\)&\(10\)&\(\displaystyle \frac2{35}<r_{10}<\frac{61}{1060}\)\\[1mm]
\(81/8\)&\(81/8\)&\(\displaystyle \frac{13}{189}<r_{81/8}<\frac{11}{159}\)\\[1mm]
\(21/2\)&\(21/2\)&\(\displaystyle \frac5{49}<r_{21/2}<\frac{38}{371}\)\\[1mm]
\(\sqrt{7073}/8\)&\(1051/100<w<1052/100\)&
\(\displaystyle \frac{757}{7357}<r_w<\frac{2903}{27878}\)\\[1mm]
\(\sqrt{114}\)&\(1067/100<w<107/10\)&
\(\displaystyle \frac{79}{679}<r_w<\frac{676}{5671}\)
\end{tabular}
\end{center}

The four gaps between consecutive rational enclosures are, in order,
\begin{equation}
 \frac{2251}{200340},\qquad
 \frac{256}{7791},\qquad
 \frac{183}{389921},\qquad
 \frac{231225}{18929162},
 \label{eq:crossing-gaps}
\end{equation}
and are all positive.  For example, the first is
\(13/189-61/1060\), and similarly for the other three.  The endpoint
margins are also explicit:
\[
 \rhozero<\frac{10}{183}<\frac2{35},
 \quad
 \frac2{35}-\frac{10}{183}=\frac{16}{6405}>0,
\]
and, since \(\rhoc>7/51\),
\[
 \frac{676}{5671}<\frac7{51}<\rhoc,
 \quad
 \frac7{51}-\frac{676}{5671}=\frac{5221}{289221}>0.
\]
We have therefore proved the complete strict order
\begin{equation}
 \rhozero<r_{10}<r_{81/8}<r_{21/2}
 <r_{\sqrt{7073}/8}<r_{\sqrt{114}}<\rhoc.
 \label{eq:moving-order}
\end{equation}
This single table replaces the ten crossing-location predicates and the
four crossing-order predicates of the executable registry.

\section{Packet B: checkpoint boundary table and exhaustiveness}
\label{sec:checkpoint}

The ball-zero input consists of the following strict rational lower
bounds:
\begin{equation}
\begin{array}{c|cccc}
n=2,\ \ell&1&2&3&5\\ \hline
q_{\ell,2}^2&(77/10)^2&9^2&(103/10)^2&(129/10)^2
\end{array}
\qquad
\begin{array}{c|cc}
n=3,\ \ell&1&2\\ \hline
q_{\ell,3}^2&(21/2)^2&(61/5)^2.
\end{array}
\label{eq:zero-anchors}
\end{equation}
Angular propagation gives the lower bound \eqref{eq:propagated-B} at all
larger angular indices.  The seven first-excluded channels have the exact
cuts shown below.  Here \(M=m(m+1)\), \(B=\frB_n(\ell_0)\),
\(Z=B-M\), \(C=(M-\ell_0(\ell_0+1))/(n^2-1)\), and the last column is
the decisive incompatibility margin \(Z-C\).

\begin{center}
\small
\renewcommand{\arraystretch}{1.22}
\begin{tabular}{c|c|c|c|c|c|c}
\(m\)&\(n\)&first excluded \(\ell_0\)&\(B\)&\(Z\)&\(C\)&\(Z-C\)\\ \hline
7&2&2&\(81\)&\(25\)&\(50/3\)&\(25/3\)\\
8&2&4&\(11409/100\)&\(4209/100\)&\(52/3\)&\(7427/300\)\\
9&2&5&\(16641/100\)&\(7641/100\)&\(20\)&\(5641/100\)\\
10&2&6&\(17841/100\)&\(6841/100\)&\(68/3\)&\(13723/300\)\\
11&2&5&\(16641/100\)&\(3441/100\)&\(34\)&\(41/100\)\\
10&3&2&\(3721/25\)&\(971/25\)&\(13\)&\(646/25\)\\
11&3&2&\(3721/25\)&\(421/25\)&\(63/4\)&\(109/100\)
\end{tabular}
\end{center}

Every table entry is reconstructible from \eqref{eq:zero-anchors}.  The
only propagated base values not already anchors are
\[
 \frac{11409}{100}=\left(\frac{103}{10}\right)^2+(4\cdot5-3\cdot4),
 \qquad
 \frac{17841}{100}=\left(\frac{129}{10}\right)^2+(6\cdot7-5\cdot6).
\]
All seven margins are positive; the smallest is \(41/100\), so none of
these exclusions relies on an equality convention.  By
Lemma~\ref{lem:cut-monotonicity}, each row excludes not merely the displayed
channel but its entire angular tail.

It remains to rule out the radial tails.  On the high strip
\(\rhoc\leq\rho<1\),
\[
 z\geq z_{\rhoc}=\pi+\frac12>
 \frac{333}{106}+\frac12=\frac{193}{53},
 \qquad
 z^2>\frac{37249}{2809}.
\]
The exact margins needed below are
\begin{equation}
 \frac{37249}{2809}-\frac{90}{8}
 =\frac{22591}{11236}>0,
 \qquad
 \frac{37249}{2809}-\frac{132}{15}
 =\frac{62649}{14045}>0.
 \label{eq:radial-margins}
\end{equation}
At \(m\leq9\), an \(n=3\) channel would require
\(8z^2<M-\ell(\ell+1)\leq90\), contradicting the first inequality.
At \(m\leq11\), every \(n\geq4\) channel would require
\(15z^2\leq(n^2-1)z^2<M-\ell(\ell+1)\leq132\), contradicting the
second inequality.  These are the two base radial exclusions.

\begin{theorem}[Exhaustive checkpoint inventories]
\label{thm:inventories}
For \(\rhoc\leq\rho<1\), the channels that can satisfy the strict proxy
\eqref{eq:strict-proxy} at the five checkpoints are contained in the
following lists:
\begin{center}
\begin{tabular}{c|l}
checkpoint&possible modes\\ \hline
\(k_7\)&first \(0{:}6\), second \(0{:}1\), no \(n\geq3\)\\
\(k_8\)&first \(0{:}7\), second \(0{:}3\), no \(n\geq3\)\\
\(k_9\)&first \(0{:}8\), second \(0{:}4\), no \(n\geq3\)\\
\(k_{10}\)&first \(0{:}9\), second \(0{:}5\), third \(0{:}1\), no \(n\geq4\)\\
\(k_{11}\)&first \(0{:}10\), second \(0{:}4\), third \(0{:}1\), no \(n\geq4\).
\end{tabular}
\end{center}
Here \(n\)-th means radial index \(n\), and \(a{:}b\) denotes the angular
indices \(a,a+1,\ldots,b\).  The word ``possible'' asserts an exhaustive
upper inventory, not that every listed channel is present for every ratio.
\end{theorem}

\begin{proof}
For \(n=1\), the product part of the strict proxy at \(k_m\) is
\[
 z^2+\ell(\ell+1)<z^2+m(m+1).
\]
Thus \(\ell<m\).  The equality channel \(\ell=m\) and all larger
indices are excluded, with the equality ownership fixed by
Lemma~\ref{lem:strict-bracket}.  This gives precisely first
\(0{:}m-1\).

For \(n=2\), apply the first five rows of the checkpoint table and then
Lemma~\ref{lem:cut-monotonicity}.  They give first-excluded angular indices
\(2,4,5,6,5\) at \(m=7,8,9,10,11\), respectively.  Hence the second
blocks end at \(1,3,4,5,4\).

For \(n=3\), the first inequality in \eqref{eq:radial-margins} removes
all channels through \(k_9\).  At \(k_{10}\) and \(k_{11}\), the last
two rows of the checkpoint table and cut monotonicity remove every
\(\ell\geq2\), leaving only \(\ell=0,1\).  Finally, the second
inequality in \eqref{eq:radial-margins} removes all \(n\geq4\) through
\(k_{11}\).  The five lists follow, without any bounded angular search.
\end{proof}

\section{Executable-row coverage contract}
\label{sec:coverage}

The former executable baseline contained 611 rows and had canonical digest
\begin{center}
\texttt{afd7e054f52aa9a3992fc8ef16d3e7551586c76bfb589053714f5f55a36792f1}.
\end{center}
The digest is included only to identify the audited baseline.  It is not a
premise of any proof above.  The exhibits above replace disjoint row groups
as follows:
\begin{center}
\small
\renewcommand{\arraystretch}{1.16}
\begin{tabularx}{\textwidth}{@{}p{0.27\textwidth}|r|>{\raggedright\arraybackslash}X@{}}
analytic exhibit&rows&exact ledger subfamilies\\ \hline
Lemma~\ref{lem:pi}&4&one enclosure-nonemptiness row, two rational
\(\pi\)-bound rows, and one reciprocal row\\
Lemma~\ref{lem:strict-bracket}&24&sixteen integer/half-integer strict
brackets, seven phase-proxy integer walls, and one optical radial equality
owner\\
Lemma~\ref{lem:block-cap}&97&twelve strict angular equality blocks, twelve
optical angular blocks, 62 high-checkpoint cap identities, one six-cap lower
tuple, and ten lower-\(d\) cap identities\\
Proposition~\ref{prop:product-deficit}&23&five general product-deficit
predicates and eighteen sampled expansions\\
Lemma~\ref{lem:cut-monotonicity}, the checkpoint boundary table, and
Theorem~\ref{thm:inventories}&117&76 repeated cut-monotonicity rows, seven
first-excluded comparisons, two radial exclusions, and 32 metadata/tail
instances\\
Lemma~\ref{lem:fixed-order} and the moving-crossing table&23&five fixed-wall
orders, one propagation identity, three endpoint comparisons, ten crossing
locations, and four crossing-order rows\\ \hline
\textbf{total}&\textbf{288}&
\end{tabularx}
\end{center}

The following family table gives the same accounting against the executable
entry points and makes the partial-family boundaries explicit.  It maps the
exact rows
replaced by this note.  ``Collapsed'' means that multiple sampled or
bookkeeping rows are consequences of one displayed identity; it does not
mean that an unprinted program enumeration is being used.

\begin{center}
\small
\renewcommand{\arraystretch}{1.18}
\begin{tabularx}{\textwidth}{@{}>{\raggedright\arraybackslash}p{0.29\textwidth}|r|r|>{\raggedright\arraybackslash}X@{}}
executable family&family rows&replaced here&analytic replacement\\ \hline
\path{verify_pi}&4&4&Lemma~\ref{lem:pi}\\
\path{verify_thresholds}&34&23&five fixed-wall orders, one propagation identity, three endpoint comparisons, and fourteen moving-crossing predicates\\
\path{verify_strict_face_conventions}&35&35&Lemmas~\ref{lem:strict-bracket} and \ref{lem:block-cap}\\
\path{verify_lower_inventories_and_payments}&44&1&single six-cap tuple bookkeeping row\\
\path{verify_checkpoint_inventories}&117&117&boundary table, radial margins, and Theorem~\ref{thm:inventories}\\
\path{verify_cap_tables_and_payments}&98&62&all block-cap arithmetic: 61 cap-table rows plus the cap-74 block identity\\
\path{verify_optical_constants}&50&36&twelve angular blocks, one radial equality owner, five general product-deficit facts, and eighteen sampled expansions\\
supplemental \(k_6\)/lower-\(d\) ledger&26&10&ten lower-\(d\) block-cap identities\\ \hline
\textbf{total}& &\textbf{288}&
\end{tabularx}
\end{center}

For the checkpoint family, the 117 rows decompose exactly as follows:
32 finite metadata/tail instances, seven first-excluded zero/channel
comparisons, two base radial exclusions, and 76 repeated cut-monotonicity
instances.  Theorem~\ref{thm:inventories} replaces all four groups.  For
the optical family, Proposition~\ref{prop:product-deficit} replaces the
18 sampled \((N,t)\) expansions by a symbolic identity valid for every
integer \(N\geq1\) and every \(0<t\leq1\).

The remaining 323 rows are outside Packets A and B.  In particular, this
note makes no claim to replace the high-event cap/payment implications,
the Bessel sign registry, seam payments, optical screen constants, or the
exact owner subtraction.  Those require the other analytic replacement
packets.  This separation prevents a universal arithmetic lemma from being
mistaken for proof of a geometric event assignment.

\end{document}
